%============================ ESCALABILIDAD (ENTREGA 3) =====================
\chapter{Optimización y escalabilidad}
\label{cap:escalabilidad}

\section{Oportunidades de optimización identificadas}
A partir de la evaluación de desempeño se identificaron mejoras concretas, resumidas en la Tabla~\ref{tab:mejoras}, clasificadas por su impacto y esfuerzo.

\begin{table}[H]
\centering
\caption{Oportunidades de optimización}
\label{tab:mejoras}
\footnotesize
\begin{tabularx}{\linewidth}{@{}X l l@{}}
\toprule
\textbf{Mejora} & \textbf{Impacto} & \textbf{Esfuerzo} \\
\midrule
Migración a base de datos relacional (SQLite/PostgreSQL) & Alto & Medio \\
Cifrado de contraseñas con \emph{hashing} (bcrypt) & Alto & Bajo \\
Limitación de tasa (\emph{rate limiting}) en el login & Medio & Bajo \\
HTTPS en el portal cautivo del ESP32 & Medio & Alto \\
Compresión de tramas WebSocket & Bajo & Bajo \\
Soporte multi-dispositivo (varios ESP32) & Alto & Medio \\
Panel de administración de invitados con filtros & Medio & Medio \\
\bottomrule
\end{tabularx}
\end{table}

\section{Plan de escalabilidad}
\subsection{Escalabilidad de datos}
El desacoplamiento del módulo \code{lib/store} (Capítulo~\ref{cap:basedatos}) permite sustituir la persistencia en archivos por un motor relacional sin modificar la lógica de negocio. Para retención prolongada y consultas analíticas, PostgreSQL con índices por marca temporal es la opción recomendada.

\subsection{Escalabilidad de dispositivos}
El modelo actual asume un dispositivo principal identificado por \code{deviceId}. Para soportar múltiples unidades, basta con particionar el estado por \code{deviceId} y extender el mapa de calor y la tabla de dispositivos con dicha dimensión, manteniendo una vista agregada por sede.

\subsection{Escalabilidad de servicio}
En caso de un número elevado de clientes del dashboard, la difusión WebSocket puede escalar horizontalmente introduciendo un \emph{message broker} (por ejemplo, Redis \emph{pub/sub}) que replique los eventos entre varias instancias del servidor tras un balanceador. Cloudflare, en el borde, aporta protección frente a picos de tráfico y denegación de servicio.

\subsection{Endurecimiento de seguridad}
Las siguientes acciones elevan la postura de seguridad para producción: almacenar contraseñas con \emph{hashing} y sal (bcrypt), aplicar limitación de intentos en el \emph{login}, rotar el token de dispositivo, y registrar auditoría de accesos. Estas mejoras son compatibles con la arquitectura actual y de bajo esfuerzo de integración.

\section{Hoja de ruta}
La Figura~\ref{fig:roadmap} propone una hoja de ruta por fases.

\begin{figure}[H]
\centering
\begin{tikzpicture}[
  ph/.style={rectangle,rounded corners=3pt,draw,thick,align=center,minimum width=3.0cm,minimum height=1.2cm,font=\small},
  arr/.style={-{Latex[length=2.4mm]},very thick}
]
\node[ph,fill=eilorcyan!12] (f1) {Fase 1\\Seguridad\\(bcrypt, rate limit)};
\node[ph,fill=eiloremerald!12,right=1.4cm of f1] (f2) {Fase 2\\Base de datos\\(PostgreSQL)};
\node[ph,fill=orange!12,right=1.4cm of f2] (f3) {Fase 3\\Multi-dispositivo\\y broker};
\node[ph,fill=gray!12,right=1.4cm of f3] (f4) {Fase 4\\SaaS multi-cliente};
\draw[arr] (f1)--(f2); \draw[arr] (f2)--(f3); \draw[arr] (f3)--(f4);
\end{tikzpicture}
\caption{Hoja de ruta de evolución del sistema.}
\label{fig:roadmap}
\end{figure}
